\documentclass[11pt]{article}

\usepackage[margin=1in]{geometry}
\usepackage{booktabs}
\usepackage{threeparttable}
\usepackage{caption}
\usepackage{float}

\captionsetup{
  font=small,
  labelfont=bf
}

\title{Accuracy and Hallucination Results}
\author{}
\date{}

\begin{document}

\maketitle

\begin{table}[H]
\centering
\caption{Accuracy results on objective QA datasets. MedQA-US is a multiple-choice benchmark, and PubMedQA is a yes/no/maybe benchmark.}
\label{tab:objective_accuracy}
\begin{threeparttable}
\small
\setlength{\tabcolsep}{9pt}
\begin{tabular}{lccccc}
\toprule
Dataset & Baseline & CoT & Meta & MedRAG (fixed-k) & MedRAG (adaptive-k) \\
\midrule
MedQA-US & 0.8174 & \textbf{0.8783} & 0.7609 & 0.8174 & 0.8174 \\
PubMedQA & 0.2240 & 0.2520 & 0.2240 & \textbf{0.6400} & \textbf{0.6400} \\
\bottomrule
\end{tabular}
\begin{tablenotes}[flushleft]
\footnotesize
\item Objective QA datasets are evaluated using exact-match accuracy.
\end{tablenotes}
\end{threeparttable}
\end{table}

\begin{table}[H]
\centering
\caption{BERTScore F1 results on subjective QA datasets. BioASQ consists of factoid questions, and MediQA contains candidate-set style subjective responses.}
\label{tab:subjective_accuracy}
\begin{threeparttable}
\small
\setlength{\tabcolsep}{9pt}
\begin{tabular}{lccccc}
\toprule
Dataset & Baseline & CoT & Meta & MedRAG (fixed-k) & MedRAG (adaptive-k) \\
\midrule
BioASQ & 0.8480 & 0.7846 & 0.8227 & 0.9049 & \textbf{0.9070} \\
MediQA & \textbf{0.8373} & 0.8203 & 0.8312 & 0.8344 & 0.8345 \\
\bottomrule
\end{tabular}
\begin{tablenotes}[flushleft]
\footnotesize
\item Subjective QA datasets are evaluated using BERTScore F1.
\end{tablenotes}
\end{threeparttable}
\end{table}

\begin{table}[H]
\centering
\caption{Hallucination rates across datasets and conditions.}
\label{tab:hallucination}
\begin{threeparttable}
\small
\setlength{\tabcolsep}{9pt}
\begin{tabular}{lccccc}
\toprule
Dataset & Baseline & CoT & Meta & MedRAG (fixed-k) & MedRAG (adaptive-k) \\
\midrule
BioASQ   & 0.0882 & \textbf{0.0147} & 0.0882 & 0.0588 & 0.0809 \\
MedQA-US & 0.1680 & \textbf{0.0960} & 0.1720 & 0.1720 & 0.1800 \\
MediQA   & 0.0240 & \textbf{0.0080} & 0.0200 & 0.0480 & 0.0600 \\
PubMedQA & 0.4200 & 0.2040 & 0.4200 & 0.1880 & \textbf{0.1760} \\
\bottomrule
\end{tabular}
\begin{tablenotes}[flushleft]
\footnotesize
\item Lower hallucination rate indicates fewer hallucinated responses.
\end{tablenotes}
\end{threeparttable}
\end{table}

\end{document}tabe